\documentclass[tikz,border=10pt]{standalone}
\usepackage[no-math]{fontspec}
\usepackage[brazil]{babel}
\usepackage{amsmath,amssymb}
\IfFileExists{C:/Windows/Fonts/Lato_Regular.ttf}{%
  \setmainfont{Lato}[
    Path = C:/Windows/Fonts/,
    UprightFont = Lato_Regular.ttf,
    BoldFont = Lato_Bold.ttf,
    ItalicFont = Lato_Italic.ttf,
    BoldItalicFont = Lato_BoldItalic.ttf
  ]
  \setsansfont{Lato}[
    Path = C:/Windows/Fonts/,
    UprightFont = Lato_Regular.ttf,
    BoldFont = Lato_Bold.ttf,
    ItalicFont = Lato_Italic.ttf,
    BoldItalicFont = Lato_BoldItalic.ttf
  ]
}{%
  \RequirePackage[default]{lato}
}
\renewcommand{\familydefault}{\sfdefault}

\usepackage{xcolor}
\usepackage{tikz}
\usepackage{pgfplots}
\pgfplotsset{compat=1.18}
\usetikzlibrary{arrows.meta,positioning,shapes.geometric,fit,backgrounds,calc,decorations.pathreplacing,patterns}

% Paleta institucional UEFS / Bioengenharia de Solos
\definecolor{bsNavy}{HTML}{1F4268}
\definecolor{bsNavyDark}{HTML}{152E4A}
\definecolor{bsGold}{HTML}{B38A3D}
\definecolor{bsGoldDark}{HTML}{8C6927}
\definecolor{bsTeal}{HTML}{0E7C7B}
\definecolor{bsForest}{HTML}{2D6A4F}
\definecolor{bsAlert}{HTML}{9E2A2B}
\definecolor{bsEarth}{HTML}{8C5A3C}
\definecolor{bsWater}{HTML}{2E86AB}
\definecolor{bsInk}{HTML}{26313D}
\definecolor{bsMuted}{HTML}{5B6573}

\definecolor{bgNavy}{HTML}{EDF2F7}
\definecolor{bgGold}{HTML}{F8F3E9}
\definecolor{bgTeal}{HTML}{EBF6F5}
\definecolor{bgForest}{HTML}{EEF5F0}
\definecolor{bgAlert}{HTML}{FDF2F2}

\tikzset{
  >=Stealth,
  bsBox/.style={
    draw=#1, fill=white, rounded corners=5pt, line width=1.15pt,
    inner sep=7pt, align=center, font=\footnotesize\sffamily, text=bsInk
  },
  bsTint/.style={
    draw=#1, fill=#1!8!white, rounded corners=5pt, line width=1.15pt,
    inner sep=7pt, align=center, font=\footnotesize\sffamily, text=bsInk
  },
  bsBadge/.style={
    fill=#1, text=white, rounded corners=9pt, inner xsep=8pt, inner ysep=3pt,
    font=\scriptsize\bfseries\sffamily
  },
  bsTitle/.style={font=\normalsize\bfseries\sffamily, text=bsNavy, align=center},
  bsSub/.style={font=\footnotesize\sffamily, text=bsMuted, align=center},
  bsArrow/.style={->, >=Stealth, line width=1.35pt},
  bsDArrow/.style={<->, >=Stealth, line width=1.35pt}
}

\tikzset{
  pill/.style={bsBadge=#1, minimum height=0pt, minimum width=0pt, text width=}
}

\begin{document}
\begin{tikzpicture}

% Faixa superior: Conectividade hidrossedimentológica (4 blocos alinhados)
\node[bsTint=bsNavy, text width=3.45cm, minimum height=2.20cm] (t1) at (0, 2.8) {%
  \tikz\node[pill=bsNavy]{ORIGEM};\\[5pt]
  \textbf{Escoamento de Vertente\\\& Estradas}\\[3pt]
  {\scriptsize Enxurrada concentrada e\\aporte de sedimentos}
};

\node[bsTint=bsGold, text width=3.45cm, minimum height=2.20cm] (t2) at (4.65, 2.8) {%
  \tikz\node[pill=bsGoldDark]{INCISÃO};\\[5pt]
  \textbf{Cabeceira / Sulcos}\\[3pt]
  {\scriptsize Concentração de energia\\e aprofundamento do canal}
};

\node[bsTint=bsTeal, text width=3.45cm, minimum height=2.20cm] (t3) at (9.30, 2.8) {%
  \tikz\node[pill=bsTeal]{CONTROLE FÍSICO};\\[5pt]
  \textbf{Série de Paliçadas\\Permeáveis}\\[3pt]
  {\scriptsize Quebra de energia e\\filtração escalonada}
};

\node[bsTint=bsForest, text width=3.45cm, minimum height=2.20cm] (t4) at (13.95, 2.8) {%
  \tikz\node[pill=bsForest]{RESULTADO};\\[5pt]
  \textbf{Retenção de Sedimentos\\\& Menos Assoreamento}\\[3pt]
  {\scriptsize Cunha deposicional estável\\e proteção a jusante}
};

\draw[bsArrow, bsNavy] (t1) -- (t2);
\draw[bsArrow, bsGoldDark] (t2) -- (t3);
\draw[bsArrow, bsTeal] (t3) -- (t4);

% Divisor conceitual bem espaçado
\draw[dashed, bsMuted!55, line width=0.9pt] (-2.05, 1.05) -- (16.00, 1.05);
\node[fill=white, inner xsep=10pt, font=\scriptsize\bfseries\sffamily, text=bsNavy]
  at (6.975, 1.05) {EXECUÇÃO EM CAMPO — SEQUÊNCIA CONSTRUTIVA EM 4 ETAPAS};

% Faixa inferior: 4 etapas construtivas
\node[bsBox=bsNavy, text width=3.45cm, minimum height=2.35cm] (b1) at (0, -0.75) {%
  \tikz\node[pill=bsNavy]{PASSO 1};\\[5pt]
  \textbf{Locação em Série}\\[4pt]
  {\scriptsize Espaçamento definido pelo\\declive e altura útil ($H$);\\crista alinhada à base superior}
};

\node[bsBox=bsEarth, text width=3.45cm, minimum height=2.35cm] (b2) at (4.65, -0.75) {%
  \tikz\node[pill=bsEarth]{PASSO 2};\\[5pt]
  \textbf{Engastamento Lateral\\e Basal}\\[3pt]
  {\scriptsize Embutimento nas paredes e\\no leito para evitar \textit{piping}\\lateral e solapamento}
};

\node[bsBox=bsTeal, text width=3.45cm, minimum height=2.35cm] (b3) at (9.30, -0.75) {%
  \tikz\node[pill=bsTeal]{PASSO 3};\\[5pt]
  \textbf{Montagem e Amarrilhos}\\[4pt]
  {\scriptsize Estacas + sacaria de ráfia\\ou bambu-madeira com\\tela filtrante e amarrilhos}
};

\node[bsBox=bsForest, text width=3.45cm, minimum height=2.35cm] (b4) at (13.95, -0.75) {%
  \tikz\node[pill=bsForest]{PASSO 4};\\[5pt]
  \textbf{Plantio Pós-Retenção}\\[4pt]
  {\scriptsize Implantação de Vetiver e\\gramíneas sobre a cunha de\\sedimentos acumulada}
};

\draw[bsArrow, bsNavy] (b1) -- (b2);
\draw[bsArrow, bsEarth] (b2) -- (b3);
\draw[bsArrow, bsTeal] (b3) -- (b4);

% Título e subtítulo ancorados acima da faixa superior
\node[bsSub] (sub) at (6.975, 4.55)
  {Integração entre fluxo de vertente, controle no talvegue e estabilização biológica pós-retenção};
\node[bsTitle, above=0.14cm of sub] (title)
  {Conectividade Hidrossedimentológica e Construção de Paliçadas em Ravinas};

\end{tikzpicture}
\end{document}
